\par\bigskip\pdfbookmark[0]{Über den Zusammenhang. Fragment V}{nexus-fragment-V}\centerline{[V.]}\medskip
Die durchschnittliche Anzahl der Klassen\begingroup\renewcommand{\thefootnote}{*)}\footnote{[Vgl. \emph{Disquiss. Arithm.}, Artikel 302; die dortige Formel weicht um eine Konstante $\delta$ von der hier im Text vorkommenden ab.]}\endgroup\ in der Nähe der negativen Determinante $-D$ ist näherungsweise
\[
=\frac{\pi\surd D}{4(1+\tfrac1{27}+\tfrac1{125}+\text{usw.})}
\]
Die tatsächliche Anzahl wird durch die folgenden Formeln ausgedrückt. Der Kürze wegen steht dabei $m$ für die durchschnittliche, $M$ für die tatsächliche Anzahl; $p,q$ bezeichnen alle ungeraden Primzahlen, die $D$ nicht teilen, und zwar ersteres die Teiler, letzteres die Nichtteiler von $\square+D$; $r$ die Primzahlen\begingroup\renewcommand{\thefootnote}{**)}\footnote{[ungerade.]}\endgroup, die $D$ teilen:

\SourceLeaf{295}
\[
\begin{aligned}
\text{I.}\quad M&=m\ \text{Produkt aus }\frac{p^3+p^2}{p^3-1}\pdot\frac{q^3-q^2}{q^3-1}\pdot\frac{r^3-r}{r^3-1}\\
\text{II.}\quad M&=\frac{\pi\surd D}4\ \text{Produkt aus }\frac{p+1}p\pdot\frac{q-1}q\pdot\frac{rr-1}{rr}\\
\text{III. NB.}\quad M&=\frac{2\surd D}\pi\ \text{Produkt aus }\frac p{p-1}\pdot\frac q{q+1}\\
\text{IV.}\quad M&=\surd\left\{\frac D2\pdot\text{Produkt aus }\frac{p+1}{p-1}\pdot\frac{q-1}{q+1}\pdot\frac{rr-1}{rr}\right\}\\
\text{V.}\quad M&=\frac{2\surd D}\pi\{1\pm\tfrac13\pm\tfrac15\pm\tfrac17\pm\tfrac19\text{ usw.}\}
\end{aligned}
\]
NB. Die Formel III wird unmittelbar aus dem Vergleich der beiden Arten abgeleitet, die darstellbaren Zahlen bis zu einer gewissen Grenze zu zählen.
\BookRule{.30\textwidth}
